\uuid{hy2g}
\titre{Fraction}
\auteur{bodin ; hanani ; mzari}
\organisation{exo7}
\licence{CC BY-NC-SA 4.0}
\datecreate{2020-01-20}
\niveau{L1}
\module{Analyse}
\chapitre{Propriétés de R}
\sousChapitre{Les rationnels}
\difficulte{1}
\theme{Rationnels, fraction}
\exo7qcmid{lille-1.reels.001}

\contenu{
\begin{question}
\qclassification{120.01}
\qtags{motcle=fraction}
Quelles sont les assertions vraies ?
\begin{answers}
    \bad{$\frac{4}{16}+\frac{4}{20} = \frac{9}{16}$}

    \good{$\frac{14}{12}+\frac{12}{14} = \frac{85}{42}$}

    \good{$\frac{36}{5} - 3 = \frac{21}{5}$}

    \good{$\frac{14}{15}/\frac{21}{35} = \frac{14}{9}$}

\end{answers}
\begin{explanations}
Pour additionner deux fractions rationnels, il faut d'abord les réduire au même dénominateur.
\end{explanations}
\end{question}
}
